% Part 1 of Day 12: Local decision-making.
%
% Topics of the syllabus: hybrid architecture patterns; which decisions must
% stay local; rules, state machines, hysteresis; safe behaviour when a part
% fails. Exercise: write the state machine for a door monitor with two
% sensors.
%
% The numbers of this part have three origins. Each frame says so in its
% source line.
% 1. A simulation of this course: tools/experiments/day12_decisions.py,
%    values in tools/data/day12/day12_decisions.json. The signals are
%    synthetic: no sensor and no model made them. The plots use the stored
%    values (seed 0); the table of the rules uses the mean of 100 seeds.
% 2. The textbook (hybrid patterns), the kit lab (the cascade of a voice
%    assistant), and the ESP32-S3 data sheet (watchdog timers).
% 3. Values of Day 11 (the round trip to a server on the internet, the
%    resend time of TCP on Linux).

\theorypart{Local decision-making}%
  {write the state machine for a door monitor with two sensors}

% ------------------------------------------------------------------ decision
\begin{frame}{A model gives a score, a program decides}
  \begingroup\centering
  \begin{tikzpicture}[font=\scriptsize, >=stealth,
      st/.style={draw=coolgray, rounded corners=2pt, fill=boxgray,
                 align=center, text width=1.75cm, minimum height=0.95cm}]
    \node[st] (s) at (0,0) {Sensor};
    \node[st, fill=cardinalred!12] (m) at (2.45,0) {Model\\ a score each
      window};
    \node[st, fill=darkcyan!20, text width=2.2cm] (d) at (5.15,0)
      {Decision layer\\ rules, states, timers};
    \node[st] (a) at (7.85,0.55) {Action};
    \node[st] (n) at (7.85,-0.55) {Message};
    \draw[->] (s) -- (m); \draw[->] (m) -- (d);
    \draw[->] (d) -- (a); \draw[->] (d) -- (n);
  \end{tikzpicture}\par\endgroup
  \vskip 0.5em
  \small
  \begin{itemize}
    \item A model gives a score for each window: ``person 0.62''. It does
          not say what the system must do.
    \item The decision layer turns the scores into actions: switch on a
          light, stop a motor, send an alert.
    \item It is ordinary code: rules, a state machine, and timers. You can
          read it, test it, and change it with no new training.
    \item The scores change from window to window. The decision layer must
          stay stable when a score jumps for one window.
  \end{itemize}
  \keypoint{A good model with a bad decision layer gives a bad system.}
\end{frame}

% ------------------------------------------------------------------ patterns
\begin{frame}{Hybrid architecture patterns}
  \footnotesize
  \renewcommand{\arraystretch}{1.18}
  \setlength{\tabcolsep}{4pt}
  \begin{tabular}{@{}L{0.19\textwidth}L{0.23\textwidth}L{0.25\textwidth}L{0.24\textwidth}@{}}
    \toprule
    \textbf{Pattern} & \textbf{Trade-off} & \textbf{Choose when} &
      \textbf{In this course} \\
    \midrule
    Train-serve split & Training cost against inference latency &
      Training needs scale that inference does not need & Train on a
      laptop, run on the XIAO (Days 3 to 5) \\
    Hierarchical processing & Local autonomy against global optimization &
      The data is more than the link carries; decisions on several time
      scales & XIAO, then Raspberry Pi, then a central computer (this
      lab) \\
    Progressive deployment & Model quality against deployment reach & The
      same model at several capability levels & One detector in
      \code{float32} and \code{int8} (Day 8) \\
    \bottomrule
  \end{tabular}
  \vskip 0.5em
  \small
  \begin{itemize}
    \item A hybrid system puts each stage on the tier where it meets its
          limits best: the time, the energy, the memory, or the data.
    \item The tiers must agree on the model version, the features, and the
          pre-processing. Else the same input gives two answers.
  \end{itemize}
  \source{\emph{Machine Learning Systems}, Volume I, chapter ``ML
    Systems'', section ``Hybrid Architectures''.}
\end{frame}

% ------------------------------------------------------------------ hierarchy
\begin{frame}{Decisions on three time scales}
  \begingroup\centering
  \begin{tikzpicture}[font=\scriptsize, >=stealth,
      tier/.style={draw=coolgray, rounded corners=2pt, align=center,
                   text width=2.3cm, minimum height=1.25cm}]
    \node[tier, fill=cardinalred!12] (x) at (0,0)
      {\textbf{XIAOML Kit}\\ a class each 0.5 s\\ decides in ms};
    \node[tier, fill=darkcyan!15] (p) at (4.0,0)
      {\textbf{Raspberry Pi}\\ rules for a room\\ decides in s};
    \node[tier, fill=boxgray] (c) at (8.0,0)
      {\textbf{Central computer}\\ logs, reports\\ minutes to days};
    \draw[->, line width=1pt] ([yshift=4pt]x.east) -- node[above] {results}
      ([yshift=4pt]p.west);
    \draw[->, line width=1pt] ([yshift=4pt]p.east) -- node[above]
      {summaries} ([yshift=4pt]c.west);
    \draw[->, darkcyan] ([yshift=-4pt]p.west) -- node[below] {commands}
      ([yshift=-4pt]x.east);
    \draw[->, darkcyan] ([yshift=-4pt]c.west) -- node[below] {settings}
      ([yshift=-4pt]p.east);
  \end{tikzpicture}\par\endgroup
  \vskip 0.5em
  \small
  \begin{itemize}
    \item Each tier decides what it can decide alone, and sends less data
          upward: a class, not the audio; a count, not the classes.
    \item A higher tier changes the rules of a lower tier: a threshold, a
          schedule, a model version. It does not take each decision.
    \item When the link to a higher tier fails, the lower tiers continue.
          This is the lab of today.
  \end{itemize}
\end{frame}

% ------------------------------------------------------------------ cascade
\begin{frame}{A cascade: a small model wakes a large one}
  \begingroup\centering
  \begin{tikzpicture}[font=\scriptsize, >=stealth,
      st/.style={draw=coolgray, rounded corners=2pt, align=center,
                 text width=2.3cm, minimum height=1.0cm}]
    \node[st, fill=boxgray] (mic) at (0,0) {Microphone\\ always on};
    \node[st, fill=cardinalred!12] (s1) at (3.2,0) {\textbf{Stage 1}\\
      keyword model on a small processor};
    \node[st, fill=darkcyan!15] (s2) at (7.0,0) {\textbf{Stage 2}\\ large
      model, only after the keyword};
    \draw[->] (mic) -- (s1);
    \draw[->] (s1) -- node[above] {keyword} (s2);
    \draw[->, coolgray] (s1) -- ++(0,-1.0) node[below] {no keyword: nothing
      leaves the device};
  \end{tikzpicture}\par\endgroup
  \vskip 0.4em
  \small
  \begin{itemize}
    \item A voice assistant listens all the time with a keyword model on a
          small processor. Only after the keyword does it send the audio
          to a large model.
    \item The small stage runs all the time, so it must use little energy.
          The large stage runs seldom, so it can be slow and expensive.
    \item The first stage also protects privacy: most audio never leaves
          the device.
    \item Errors of stage 1 set the behaviour: a missed keyword ignores the
          user; a false keyword sends audio for nothing.
  \end{itemize}
  \source{Kit lab ``Keyword Spotting'' of \emph{Machine Learning Systems}.}
\end{frame}

% ------------------------------------------------------------------ local
\begin{frame}{Which decisions must stay local}
  \footnotesize
  \renewcommand{\arraystretch}{1.18}
  \begin{tabular}{@{}L{0.20\textwidth}L{0.36\textwidth}L{0.36\textwidth}@{}}
    \toprule
    \textbf{Reason} & \textbf{Question} & \textbf{Example} \\
    \midrule
    Deadline & Is the deadline shorter than the round trip to the server,
      with resends? & Stop a motor within 50 ms \\
    No network & Must the decision work when the link fails? & Open a
      door for a person \\
    Safety & Can a late or a missing decision hurt somebody? & Switch off
      a heater that is too hot \\
    Privacy & Does the decision need data that must not leave? & Count
      the people in a room from the video \\
    Volume & Is the input much larger than the result? & Classify 16 000
      audio samples each second \\
    \bottomrule
  \end{tabular}
  \vskip 0.5em
  \small
  \begin{itemize}
    \item One ``yes'' is enough: the decision stays on the device.
    \item On Day 11, a server on the internet answered in 69 ms. A lost
          TCP segment waits at least 200 ms on Linux. A deadline of 50 ms
          cannot use the internet.
    \item The cloud keeps the decisions that need much data and no fast
          answer: reports, retraining, new thresholds.
  \end{itemize}
  \source{Day 11, Part 1.}
\end{frame}

% ------------------------------------------------------------------ rules
\begin{frame}[fragile]{Rules on the output of a model}
\begin{lstlisting}[style=python]
def decide(scores):                   # scores: {"person": 0.62, ...}
    label = max(scores, key=scores.get)
    if scores[label] < 0.5:           # not sure: no action
        return "none"
    if label == "fire":               # a dangerous class wins
        return "alarm"
    if label == "person":
        return "light_on"
    return "none"
\end{lstlisting}
  \small
  \begin{itemize}
    \item \textbf{A threshold}: act only when the score is high enough.
          The threshold sets the balance of missed events and false
          events.
    \item \textbf{A priority}: a dangerous class wins over a normal class,
          also with a lower score.
    \item \textbf{A default}: ``no action'' when the model is not sure, or
          when the class is unknown.
    \item A rule on one window reacts to each noisy window. The next frames
          show the effect and two fixes: hysteresis and filters in time.
  \end{itemize}
\end{frame}

% ------------------------------------------------------------------ threshold sim
\begin{frame}{A noisy reading near one threshold}
  \begingroup\centering
  \begin{tikzpicture}[font=\scriptsize]
    \draw[->, coolgray] (0,0) -- (9.3,0) node[right] {s};
    \draw[->, coolgray] (0,0) -- (0,3.3);
    \node[coolgray, anchor=west] at (0.1,3.3) {degrees};
    \foreach \t in {0,100,...,600} {
      \draw[coolgray] ({\t*0.015},0) -- ({\t*0.015},-0.08)
        node[below] {\t};
    }
    \foreach \v in {24,25,26} {
      \draw[coolgray] (0,{(\v-23.5)*1.1}) -- (-0.08,{(\v-23.5)*1.1})
        node[left] {\v};
    }
    \draw[cardinalred, dashed] (0,1.65) -- (9.0,1.65) node[right]
      {25.0};
    \draw[darkcyan, line width=0.6pt] plot coordinates {(0.000,0.571) (0.060,0.491) (0.120,0.493) (0.180,0.254) (0.240,0.577) (0.300,0.675)
      (0.360,0.876) (0.420,0.680) (0.480,0.758) (0.540,0.706) (0.600,0.636) (0.660,0.916)
      (0.720,1.199) (0.780,0.931) (0.840,1.032) (0.900,0.919) (0.960,1.074) (1.020,1.153)
      (1.080,1.244) (1.140,0.998) (1.200,1.168) (1.260,0.955) (1.320,1.070) (1.380,1.153)
      (1.440,1.280) (1.500,1.366) (1.560,1.456) (1.620,0.970) (1.680,1.358) (1.740,1.635)
      (1.800,1.560) (1.860,1.437) (1.920,1.395) (1.980,1.577) (2.040,1.779) (2.100,1.632)
      (2.160,1.659) (2.220,1.503) (2.280,1.525) (2.340,1.744) (2.400,1.808) (2.460,2.053)
      (2.520,2.054) (2.580,1.962) (2.640,1.597) (2.700,1.836) (2.760,1.827) (2.820,1.987)
      (2.880,2.137) (2.940,1.775) (3.000,1.907) (3.060,2.179) (3.120,1.998) (3.180,1.975)
      (3.240,2.133) (3.300,2.150) (3.360,2.110) (3.420,2.144) (3.480,2.450) (3.540,1.989)
      (3.600,2.525) (3.660,2.630) (3.720,2.316) (3.780,2.574) (3.840,2.318) (3.900,2.341)
      (3.960,2.485) (4.020,2.281) (4.080,2.416) (4.140,2.579) (4.200,2.677) (4.260,2.148)
      (4.320,2.556) (4.380,2.797) (4.440,2.673) (4.500,2.948) (4.560,2.793) (4.620,2.759)
      (4.680,2.952) (4.740,2.617) (4.800,2.807) (4.860,2.908) (4.920,2.761) (4.980,2.857)
      (5.040,2.803) (5.100,2.578) (5.160,2.709) (5.220,3.029) (5.280,2.685) (5.340,2.690)
      (5.400,2.523) (5.460,2.615) (5.520,2.624) (5.580,2.871) (5.640,2.905) (5.700,2.571)
      (5.760,2.619) (5.820,2.584) (5.880,2.716) (5.940,3.013) (6.000,2.691) (6.060,2.529)
      (6.120,2.686) (6.180,2.697) (6.240,2.677) (6.300,2.505) (6.360,3.098) (6.420,2.928)
      (6.480,2.907) (6.540,2.502) (6.600,2.873) (6.660,2.928) (6.720,2.842) (6.780,2.890)
      (6.840,2.603) (6.900,2.535) (6.960,2.926) (7.020,2.700) (7.080,2.718) (7.140,2.685)
      (7.200,2.840) (7.260,2.633) (7.320,2.973) (7.380,2.786) (7.440,2.715) (7.500,2.963)
      (7.560,2.953) (7.620,2.869) (7.680,2.660) (7.740,3.068) (7.800,2.806) (7.860,2.872)
      (7.920,2.733) (7.980,2.695) (8.040,2.489) (8.100,2.704) (8.160,2.883) (8.220,2.892)
      (8.280,2.809) (8.340,2.697) (8.400,2.609) (8.460,2.723) (8.520,2.926) (8.580,2.927)
      (8.640,2.718) (8.700,2.970) (8.760,2.732) (8.820,2.579) (8.880,2.755) (8.940,2.579)};
  \end{tikzpicture}\par\endgroup
  \vskip 0.3em
  \small
  \begin{itemize}
    \item A temperature rises from 24 to 26 degrees in 300 s. One reading
          each second, with noise of 0.15 degrees. A fan runs above 25.0
          degrees.
    \item The true value crosses 25.0 degrees at 150 s. The noisy reading
          crosses the line many times near this point.
    \item With one threshold, the fan switches 25 times in 600 s:
          each switch wears the relay and the motor.
  \end{itemize}
  \source{Simulation of this course (synthetic signal, seed 0). One point
    each 4 s is shown.}
\end{frame}

% ------------------------------------------------------------------ hysteresis
\begin{frame}{Hysteresis: two thresholds}
  \begin{columns}[c]
    \begin{column}{0.44\textwidth}
      \centering
      \begin{tikzpicture}[font=\scriptsize, >=stealth, x=0.92cm, y=1.2cm]
        \draw[->, coolgray] (0,0) -- (3.3,0) node[right] {reading};
        \draw[->, coolgray] (0,0) -- (0,1.5) node[above] {fan};
        \node[left, coolgray] at (0,0.2) {off};
        \node[left, coolgray] at (0,1.1) {on};
        \draw[cardinalred, line width=1pt] (0.2,0.2) -- (2.0,0.2) -- (2.0,1.1)
          -- (3.1,1.1);
        \draw[darkcyan, line width=1pt] (3.1,1.1) -- (1.2,1.1) -- (1.2,0.2);
        \draw[->, cardinalred] (1.5,0.35) -- (1.8,0.35);
        \draw[->, darkcyan] (1.8,0.95) -- (1.5,0.95);
        \draw[coolgray, dotted] (1.2,0) -- (1.2,0.2);
        \draw[coolgray, dotted] (2.0,0) -- (2.0,0.2);
        \node[below] at (1.2,0) {off below};
        \node[below] at (2.0,-0.25) {on above};
      \end{tikzpicture}
    \end{column}
    \begin{column}{0.54\textwidth}
      \footnotesize
      \renewcommand{\arraystretch}{1.15}
      \setlength{\tabcolsep}{4pt}
      \begin{tabular}{@{}L{1.9cm}rr@{}}
        \toprule
        \textbf{Rule} & \textbf{Switches} & \textbf{First on} \\
        \midrule
        Threshold 25.0 & 25 & 133 s \\
        Band $\pm 0.2$      & 1 & 164 s \\
        Band $\pm 0.4$      & 1 & 187 s \\
        \bottomrule
      \end{tabular}
      \vskip 0.2em
      {\scriptsize The same 600 readings. The true value crosses 25.0 at
      150 s.\par}
    \end{column}
  \end{columns}
  \vskip 0.4em
  \small
  \begin{itemize}
    \item Switch on above the upper threshold, and off only below the lower
          one. Noise inside the band changes nothing.
    \item Select the band wider than the noise: here about 2 to 3 times the
          standard deviation of 0.15 degrees.
    \item The cost is delay: the band of $\pm 0.4$ switched on 37 s
          after the true crossing. One threshold switched on 17 s
          before it, because of the noise.
  \end{itemize}
  \source{Simulation of this course (synthetic signal, seed 0).}
\end{frame}

% ------------------------------------------------------------------ scores sim
\begin{frame}{Noise in time: the scores of a model}
  \begingroup\centering
  \begin{tikzpicture}[font=\scriptsize]
    \fill[treegreen!15] ({20*0.135},0) rectangle ({35*0.135},2.2);
    \node[treegreen!60!black] at ({27.5*0.135},2.35) {person present};
    \draw[->, coolgray] (0,0) -- (8.4,0) node[right] {s};
    \draw[->, coolgray] (0,0) -- (0,2.5);
    \node[coolgray, anchor=west] at (0.1,2.6) {score};
    \foreach \t in {0,10,...,60} {
      \draw[coolgray] ({\t*0.135},0) -- ({\t*0.135},-0.08) node[below] {\t};
    }
    \foreach \v in {0.5,1.0} {
      \draw[coolgray] (0,{\v*2}) -- (-0.08,{\v*2}) node[left] {\v};
    }
    \draw[cardinalred, dashed] (0,1.0) -- (8.1,1.0);
    \draw[darkcyan, line width=0.5pt] plot coordinates {(0.000,0.122) (0.054,0.466) (0.108,0.078) (0.162,0.354) (0.216,0.000) (0.270,0.836)
      (0.324,0.234) (0.378,0.324) (0.432,0.698) (0.486,0.816) (0.540,0.268) (0.594,0.254)
      (0.648,0.374) (0.702,0.728) (0.756,0.580) (0.810,0.256) (0.864,0.220) (0.918,0.158)
      (0.972,0.480) (1.026,0.470) (1.080,0.122) (1.134,0.604) (1.188,0.516) (1.242,0.142)
      (1.296,0.328) (1.350,0.704) (1.404,0.434) (1.458,0.372) (1.512,0.700) (1.566,0.530)
      (1.620,0.506) (1.674,0.686) (1.728,0.844) (1.782,0.308) (1.836,0.048) (1.890,0.000)
      (1.944,0.256) (1.998,0.582) (2.052,0.730) (2.106,0.538) (2.160,0.222) (2.214,0.908)
      (2.268,0.218) (2.322,0.000) (2.376,0.544) (2.430,0.110) (2.484,0.358) (2.538,0.234)
      (2.592,0.146) (2.646,1.682) (2.700,1.300) (2.754,1.742) (2.808,1.358) (2.862,1.264)
      (2.916,1.498) (2.970,1.476) (3.024,1.350) (3.078,1.358) (3.132,1.862) (3.186,0.970)
      (3.240,1.890) (3.294,2.000) (3.348,1.402) (3.402,1.820) (3.456,1.300) (3.510,1.288)
      (3.564,1.498) (3.618,1.076) (3.672,1.266) (3.726,1.510) (3.780,1.636) (3.834,0.620)
      (3.888,1.308) (3.942,1.692) (3.996,1.412) (4.050,1.860) (4.104,1.578) (4.158,1.516)
      (4.212,1.868) (4.266,1.258) (4.320,1.604) (4.374,1.788) (4.428,1.520) (4.482,1.694)
      (4.536,1.596) (4.590,1.188) (4.644,1.426) (4.698,2.000) (4.752,1.436) (4.806,0.180)
      (4.860,0.000) (4.914,0.034) (4.968,0.122) (5.022,0.000) (5.076,0.382) (5.130,0.576)
      (5.184,0.418) (5.238,0.496) (5.292,0.294) (5.346,0.270) (5.400,0.684) (5.454,0.072)
      (5.508,0.270) (5.562,0.392) (5.616,0.658) (5.670,0.116) (5.724,0.516) (5.778,0.678)
      (5.832,0.188) (5.886,0.492) (5.940,0.662) (5.994,0.282) (6.048,0.858) (6.102,0.358)
      (6.156,0.120) (6.210,1.792) (6.264,0.862) (6.318,0.420) (6.372,0.380) (6.426,0.062)
      (6.480,0.444) (6.534,0.148) (6.588,0.320) (6.642,0.540) (6.696,0.216) (6.750,0.436)
      (6.804,0.482) (6.858,0.638) (6.912,0.450) (6.966,0.154) (7.020,0.250) (7.074,0.562)
      (7.128,0.598) (7.182,0.348) (7.236,0.208) (7.290,0.264) (7.344,0.630) (7.398,0.628)
      (7.452,0.150) (7.506,0.258) (7.560,0.000) (7.614,0.626) (7.668,0.582) (7.722,0.832)
      (7.776,0.376) (7.830,0.254) (7.884,0.182) (7.938,0.762) (7.992,0.202) (8.046,0.578)};
    \node[anchor=west] at (0,-0.75) {one window};
    \node[anchor=west] at (0,-1.15) {mean 3, hysteresis};
    \fill[cardinalred!70] (2.6460,-0.85) rectangle (2.6595,-0.65);
    \fill[cardinalred!70] (2.7000,-0.85) rectangle (2.9835,-0.65);
    \fill[cardinalred!70] (2.9970,-0.85) rectangle (3.1860,-0.65);
    \fill[cardinalred!70] (3.1995,-0.85) rectangle (3.2130,-0.65);
    \fill[cardinalred!70] (3.2265,-0.85) rectangle (3.8340,-0.65);
    \fill[cardinalred!70] (3.8475,-0.85) rectangle (4.0905,-0.65);
    \fill[cardinalred!70] (4.1040,-0.85) rectangle (4.5765,-0.65);
    \fill[cardinalred!70] (4.5900,-0.85) rectangle (4.7250,-0.65);
    \fill[cardinalred!70] (4.7520,-0.85) rectangle (4.7655,-0.65);
    \fill[cardinalred!70] (5.6565,-0.85) rectangle (5.6700,-0.65);
    \fill[cardinalred!70] (6.0885,-0.85) rectangle (6.1020,-0.65);
    \fill[cardinalred!70] (6.2100,-0.85) rectangle (6.2235,-0.65);
    \fill[cardinalred!70] (6.6690,-0.85) rectangle (6.6825,-0.65);
    \fill[cardinalred!70] (7.0875,-0.85) rectangle (7.1010,-0.65);
    \fill[cardinalred!70] (7.9785,-0.85) rectangle (7.9920,-0.65);
    \fill[darkcyan] (2.7270,-1.25) rectangle (4.7385,-1.05);
  \end{tikzpicture}\par\endgroup
  \vskip 0.3em
  \small
  \begin{itemize}
    \item A ``person'' score, 10 windows each second. A person is present
          from 20 s to 35 s. With no person, 6 single windows have a false
          high score.
    \item The bars show when a rule says ``person''. The rule on one
          window says it 15 times in this run, 8 times
          with no person.
  \end{itemize}
  \source{Simulation of this course (synthetic scores, seed 0). The line
    shows one window each 0.4 s; the bars use all windows.}
\end{frame}

% ------------------------------------------------------------------ rules table
\begin{frame}{Simulated: five rules on the same scores}
  \footnotesize
  \renewcommand{\arraystretch}{1.15}
  \begin{tabular}{@{}L{0.40\textwidth}rrrr@{}}
    \toprule
    \textbf{Rule} & \textbf{Switch-on} & \textbf{False} & \textbf{Delay on}
      & \textbf{Delay off} \\
    \midrule
    One window above 0.5 & 16.11 & 8.60 & 0.04 s & 0.10 s \\
    Mean of 3 above 0.5 & 2.32 & 1.09 & 0.11 s & 0.09 s \\
    3 of 5 above 0.5 & 1.07 & 0.00 & 0.21 s & 0.20 s \\
    Hysteresis, one window & 8.50 & 6.11 & 0.17 s & 0.01 s \\
    Hysteresis, mean of 3 & 1.06 & 0.06 & 0.17 s & 0.14 s \\
    \bottomrule
  \end{tabular}
  \vskip 0.1em
  {\scriptsize Mean of 100 runs (seeds 0 to 99). Hysteresis: on above 0.6,
  off below 0.4. The correct result is 1 switch-on and 0 false. A delay is
  the time after the true start (20 s) or end (35 s).\par}
  \vskip 0.4em
  \small
  \begin{itemize}
    \item Hysteresis alone cuts the switch-on events inside the event from
          7.51 to 2.39. It does not stop a single false window: such a
          window is above both thresholds.
    \item A filter in time (the mean of 3, or 3 of 5) stops single false
          windows. ``3 of 5'' had no false switch-on in 100 runs, for a
          delay of about 0.2 s.
  \end{itemize}
  \keypoint{Filter the scores in time, and use two thresholds. A delay of
    some windows is the price.}
  \source{Simulation of this course (synthetic scores).}
\end{frame}

% ------------------------------------------------------------------ state machine
\begin{frame}{State machines}
  \begingroup\centering
  \begin{tikzpicture}[font=\scriptsize, >=stealth,
      s/.style={draw=coolgray, circle, fill=boxgray, minimum size=1.15cm,
                align=center}]
    \node[s] (off) at (0,0) {OFF};
    \node[s, fill=darkcyan!20] (on) at (4.2,0) {ON};
    \node[s, fill=cardinalred!20] (f) at (8.2,0) {FAULT};
    \draw[->] (off) to[bend left=20] node[above] {person} (on);
    \draw[->] (on) to[bend left=20] node[below] {30 s with no person}
      (off);
    \draw[->] (on) -- node[above] {no data for 10 s} (f);
    \draw[->] (f.south) to[bend left=25] node[below] {data again}
      (off.south);
    \draw[->] (on) to[loop above, looseness=5] node[above] {person: restart
      the timer} (on);
  \end{tikzpicture}\par\endgroup
  \vskip 0.3em
  \small
  \begin{itemize}
    \item A \textbf{state} is what the system does now. An \textbf{event}
          is an input: a class, a timer, a lost sensor. A
          \textbf{transition} goes from one state to another for one
          event.
    \item Each state has a fixed action: the light is on in ON and off in
          OFF and in FAULT.
    \item Write each state and each event in a table. Each pair needs an
          answer, also ``stay''. A missing pair is a bug.
  \end{itemize}
\end{frame}

\begin{frame}[fragile]{A state machine in code}
\begin{lstlisting}[style=python]
import time
state, last_person, last_data = "OFF", 0.0, time.monotonic()
def step(person, now):                  # called for each new window
    global state, last_person, last_data
    if person is None:                  # no input in this window
        if now - last_data > 10:
            state = "FAULT"
        return state
    last_data = now
    if person:
        last_person = now
    if state == "FAULT":
        state = "OFF"
    if state == "OFF" and person:
        state = "ON"
    elif state == "ON" and now - last_person > 30:
        state = "OFF"
    return state
\end{lstlisting}
  \small
  \begin{itemize}
    \item The function never waits: no \code{sleep}. It compares times, so
          it can run next to the model and the network code.
    \item Use a monotonic clock (\code{time.monotonic()}, \code{millis()}).
          The clock of the day can jump.
  \end{itemize}
\end{frame}

% ------------------------------------------------------------------ failures
\begin{frame}{Safe behaviour when a part fails}
  \footnotesize
  \renewcommand{\arraystretch}{1.18}
  \begin{tabular}{@{}L{0.22\textwidth}L{0.34\textwidth}L{0.36\textwidth}@{}}
    \toprule
    \textbf{Failure} & \textbf{How to detect it} & \textbf{Response} \\
    \midrule
    Sensor stops & No reading for a fixed time & Safe state, a fault
      message \\
    Sensor stuck & The same value for too long, or a value out of range &
      Ignore the sensor, a fault message \\
    Model unsure & Low scores for many windows & No action, or ask a person
      \\
    Network lost & No answer, a lost connection & Continue the local
      decisions, keep the messages \\
    Device crash & No heartbeat; the watchdog timer expires & Restart in the
      safe state \\
    Power lost & The device restarts & Start in the safe state, read the
      saved settings \\
    \bottomrule
  \end{tabular}
  \vskip 0.5em
  \small
  \begin{itemize}
    \item Each input needs a timeout. A missing value is an event, not a
          zero.
    \item Plan the response before the failure, and test it: pull a cable,
          stop the broker, switch the power off.
  \end{itemize}
\end{frame}

\begin{frame}{The safe state depends on the system}
  \footnotesize
  \renewcommand{\arraystretch}{1.2}
  \begin{tabular}{@{}L{0.30\textwidth}L{0.28\textwidth}L{0.34\textwidth}@{}}
    \toprule
    \textbf{System} & \textbf{Safe state} & \textbf{Reason} \\
    \midrule
    Heater & Off & A heater that is on with no control can start a fire \\
    Conveyor belt & Stopped & A moving belt with no control can hurt a
      person \\
    Door of a fire exit & Unlocked & People must leave the building \\
    Door of a store room & Locked & Goods must stay protected \\
    Ventilation of a lab & On & Gas must leave the room \\
    \bottomrule
  \end{tabular}
  \vskip 0.5em
  \small
  \begin{itemize}
    \item ``Off'' is not always safe. The safe state comes from the risk of
          the system, not from the code.
    \item A safety function that must work in all cases needs a hardware
          path too: a fuse, a thermal switch, an emergency stop. A model
          and a program are not enough.
  \end{itemize}
  \keypoint{Write the safe state of each output before you write the
    rules.}
\end{frame}

% ------------------------------------------------------------------ watchdog
\begin{frame}{Watchdogs and heartbeats}
  \begingroup\centering
  \begin{tikzpicture}[font=\scriptsize, >=stealth]
    \draw[->, coolgray] (0,0) -- (9.3,0) node[right] {time};
    \foreach \x in {0.5,1.5,2.5,3.5} {
      \draw[darkcyan, line width=1pt] (\x,0) -- (\x,0.5);
    }
    \node[darkcyan, anchor=south] at (2.0,0.55) {the program resets the
      timer};
    \draw[cardinalred, dashed] (3.5,0.9) -- (7.0,0.9);
    \node[cardinalred, anchor=south] at (5.25,0.95) {the program hangs};
    \draw[->, cardinalred, line width=1pt] (7.0,1.2) -- (7.0,0.05);
    \node[cardinalred, anchor=west] at (7.1,0.6) {timeout: reset};
  \end{tikzpicture}\par\endgroup
  \vskip 0.4em
  \small
  \begin{itemize}
    \item A \textbf{watchdog timer} is a counter in the chip. The program
          resets it in each loop. If the program hangs, the counter
          expires and resets the chip.
    \item The ESP32-S3 has three watchdog timers: one in each of the two
          timer groups and one in the RTC module. Each one can reset the
          processor or the system.
    \item A \textbf{heartbeat} is the same idea between devices: a short
          message each few seconds. No heartbeat for a fixed time means
          that the device is gone.
    \item MQTT has a built-in form of this: the keep-alive and the last
          will message (Part 2).
  \end{itemize}
  \source{ESP32-S3 data sheet, version 2.2, section ``Watchdog Timers''.}
\end{frame}

% ------------------------------------------------------------------ test
\begin{frame}{Test the decision layer}
  \small
  \begin{itemize}
    \item \textbf{Replay}: record the scores of a real run, then feed the
          same file to the decision layer. Each change of a rule then has a
          fair comparison, as in the simulation of this part.
    \item \textbf{One test for each transition}: a table of the state, the
          event, and the next state. Each row is one small test.
    \item \textbf{Fault injection}: stop the sensor, stop the broker, send
          a value out of range, restart the device. Check the safe state.
    \item \textbf{Count the switches}: in a test of one hour, how often did
          the output change? Many changes show a missing filter or a
          missing hysteresis.
  \end{itemize}
  \vskip 0.3em
  \footnotesize
  \renewcommand{\arraystretch}{1.12}
  \begin{tabular}{@{}lllL{0.34\textwidth}@{}}
    \toprule
    \textbf{State} & \textbf{Event} & \textbf{Next state} & \textbf{Test} \\
    \midrule
    OFF   & person         & ON    & one window with a person \\
    ON    & 30 s, no person & OFF  & a run of 31 s with no person \\
    ON    & no data, 10 s  & FAULT & stop the input for 11 s \\
    FAULT & data again     & OFF   & start the input again \\
    \bottomrule
  \end{tabular}
\end{frame}

% -------------------------------------------------------------------- exercise
\exerciseframe{a state machine for a door}{%
  \small
  A door monitor has two sensors: a contact \textbf{D} (door open or
  closed), and a person detector \textbf{P} (a camera model: a person near
  the door, or not). It must do this:
  \begin{enumerate}
    \item Door closed: nothing to do.
    \item Door open, and a person near it: normal.
    \item Door open with no person for 30 s: send one alert.
    \item The door closes: back to normal, and the alert ends.
    \item A sensor sends no value for 10 s: send a fault message.
  \end{enumerate}
  Write the states, the events, and the transitions as a table. Give the
  action of each state. Which input needs a filter in time, and which
  one?
  \vskip 0.3em
  Work with your neighbour for 8 minutes.
}

\begin{frame}{Answer: a state machine for a door}
  \footnotesize
  \renewcommand{\arraystretch}{1.1}
  \begin{tabular}{@{}llL{0.30\textwidth}@{}}
    \toprule
    \textbf{State, event} & \textbf{Next state} & \textbf{Action} \\
    \midrule
    CLOSED, D open and P & OPEN\_PERSON & none \\
    CLOSED, D open, no P & OPEN\_ALONE & start the 30 s timer \\
    OPEN\_PERSON, no P & OPEN\_ALONE & start the 30 s timer \\
    OPEN\_ALONE, P & OPEN\_PERSON & stop the timer \\
    OPEN\_ALONE, timer ends & ALERT & send one alert \\
    any open state or ALERT, D closed & CLOSED & stop the timer; end the
      alert \\
    any state, no value for 10 s & FAULT & send a fault message \\
    FAULT, both sensors send again & from D and P & the state of the
      inputs \\
    \bottomrule
  \end{tabular}
  \vskip 0.4em
  \small
  \begin{itemize}
    \item P comes from a model: filter it in time (``3 of 5'' windows), or
          one false window starts or stops the timer.
    \item D is a contact: a short debounce of some ms is enough.
    \item ALERT stays when a person comes: the door was open and alone for
          30 s. Only the closed door ends it.
  \end{itemize}
\end{frame}
